\uuid{kR9v}
\titre{Interprétation de filtres classiques}
\niveau{L2}
\module{Informatique / Signal}
\chapitre{Réseaux de neurones}
\sousChapitre{Filtres de convolution 2D}
\theme{Lissage, Sobel, Laplacien, détection de contours}
\auteur{}
\datecreate{2026-03-23}
\organisation{}
\difficulte{2}

\contenu{
\texte{
  On donne les quatre noyaux suivants :
  \[
  K_A = \frac{1}{9}\begin{pmatrix}1&1&1\\1&1&1\\1&1&1\end{pmatrix},\quad
  K_B = \begin{pmatrix}-1&-1&-1\\-1&8&-1\\-1&-1&-1\end{pmatrix},
  \]
  \[
  K_C = \begin{pmatrix}-1&0&+1\\-2&0&+2\\-1&0&+1\end{pmatrix},\quad
  K_D = \begin{pmatrix}0&1&0\\1&-4&1\\0&1&0\end{pmatrix}.
  \]
}
\begin{enumerate}
\item
  \question{Associer chaque noyau à l'un des rôles suivants : \textbf{lissage}, \textbf{détection de contours tous azimuts}, \textbf{détection de contours verticaux (Sobel X)}, \textbf{Laplacien discret}.}
  \reponse{
    \begin{itemize}
      \item $K_A$ : \textbf{lissage} (filtre moyenneur uniforme).
      \item $K_B$ : \textbf{détection de contours tous azimuts} (filtre laplacien centré sur 8, réagit dans toutes les directions).
      \item $K_C$ : \textbf{détection de contours verticaux} (filtre de Sobel X, sensible aux variations horizontales d'intensité).
      \item $K_D$ : \textbf{Laplacien discret} (approximation de $\nabla^2 I$ avec les 4-voisins).
    \end{itemize}
  }

\item
  \question{Pour $K_A$ : calculer la somme de ses coefficients. Qu'implique ce résultat sur les zones uniformes de l'image ?}
  \reponse{
    $\sum K_A = 9 \times \frac{1}{9} = 1$.
    Sur une zone uniforme (tous les pixels égaux à une constante $c$), la convolution
    donne $c \times 1 = c$ : l'intensité est préservée. Le filtre ne crée ni n'annule de
    luminosité ; c'est un filtre \textbf{passe-bas normalisé}.
  }

\item
  \question{Pour $K_B$ : calculer la somme de ses coefficients. Qu'implique ce résultat sur les zones uniformes ?}
  \reponse{
    $\sum K_B = 8 + 8\times(-1) = 8 - 8 = 0$.
    Sur une zone uniforme, la convolution donne $0$ : les zones plates apparaissent
    en noir dans la carte de sortie. Seules les transitions d'intensité produisent
    une réponse non nulle ; c'est caractéristique d'un filtre \textbf{passe-haut} ou
    détecteur de contours.
  }

\item
  \question{Pour $K_C$ : si on l'applique à une image dont toutes les colonnes sont identiques, que vaut la sortie ? Pourquoi ?}
  \reponse{
    Si toutes les colonnes sont identiques, l'intensité ne varie pas horizontalement.
    Or $K_C$ calcule $-I(\cdot, j-1) + I(\cdot, j+1)$ (pondéré), ce qui est une
    différence entre colonne de droite et colonne de gauche. Cette différence est nulle
    si toutes les colonnes sont égales. La sortie est donc entièrement \textbf{nulle}.
    $K_C$ ne réagit qu'aux variations horizontales d'intensité (i.e. aux contours verticaux).
  }

\item
  \question{Calculer $K_B + 9 K_A$ (addition élément par élément). Interpréter le résultat.}
  \reponse{
    \[
      K_B + 9 K_A =
      \begin{pmatrix}-1&-1&-1\\-1&8&-1\\-1&-1&-1\end{pmatrix}
      +
      \begin{pmatrix}1&1&1\\1&1&1\\1&1&1\end{pmatrix}
      =
      \begin{pmatrix}0&0&0\\0&9&0\\0&0&0\end{pmatrix}.
    \]
    On obtient $9$ fois le filtre identité. Cela illustre que $K_B$ peut s'interpréter
    comme \og identité (mise à l'échelle) $-$ filtre moyenneur \fg{} : il soustrait
    la moyenne locale pour ne conserver que les variations.
  }
\end{enumerate}
}
